\documentclass[review,12pt]{elsarticle}

\usepackage{graphicx}
\usepackage{amsmath}
\usepackage{amssymb}
\usepackage{booktabs}
\usepackage{hyperref}
\usepackage{xcolor}
\usepackage{lineno}
\usepackage{setspace}
\usepackage{array}
\usepackage{multirow}
\usepackage{geometry}
\usepackage{float}
\usepackage{caption}
\usepackage{subcaption}
\usepackage{siunitx}

\geometry{margin=2.5cm}
\modulolinenumbers[5]

\journal{Agricultural Water Management}

\begin{document}

\begin{frontmatter}

\title{Decoupling Canopy Thermal Anomalies from Root-Zone Water Depletion at Sub-Field Scales Using Physics-Informed Satellite Telemetry: A Pilot Study at Russell Ranch}

\author[1]{Umer Tanveer\corref{cor1}}
\ead{umertanveer@example.com}
\cortext[cor1]{Corresponding author.}

\author[2]{Jamal Ahmed}
\author[1]{Ali Hassan}

\address[1]{Department of Computer Science, University of Lahore, Pakistan}
\address[2]{Department of Agricultural Engineering, NUST, Islamabad, Pakistan}

\begin{abstract}
Conventional crop water stress monitoring systems treat atmospheric demand and root-zone soil moisture as interchangeable drivers of plant stress. This conflation results in premature or erroneous irrigation triggers when canopy surface temperatures rise due to hot, dry boundary-layer air (atmospheric demand) even while root-zone soil moisture remains within the readily available water (RAW) threshold. We term this phenomenon \textit{decoupled stress}: a state in which the leaf-level thermal anomaly ($\Delta T = T_s - T_{air}$) signals stress while root-zone depletion ($D_r$) remains below the deficit triggering threshold. Using a 21-day, 119,553-record sub-field telemetry log generated by the AquaVolt-AI Physics-Informed Machine Learning (PIML) pipeline at the Russell Ranch Sustainable Agriculture Facility (UC Davis, California), we decompose crop stress events into (i) root-zone driven, (ii) atmospherically driven, and (iii) decoupled events. Across three active crop fields—corn (\textit{Zea mays}), alfalfa (\textit{Medicago sativa}), and processing tomato (\textit{Solanum lycopersicum})—and an 8$\times$8 sector spatial grid at 10-m resolution, we quantify the prevalence, spatial distribution, and crop-type dependency of decoupled stress. Results demonstrate that $31.4\pm4.7$\% of all high-thermal-anomaly hours across the monitoring period correspond to decoupled conditions in which soil moisture remains adequate. Corn exhibited the highest decoupled stress rate ($38.2\%$), followed by tomato ($29.1\%$) and alfalfa ($21.4\%$). The relationship between vapor pressure deficit (VPD) and the onset of decoupled stress followed a piecewise threshold at VPD $> 1.5$ kPa, consistent with stomatal closure mechanisms documented in the literature. These findings demonstrate that irrigation scheduling systems relying solely on thermal canopy indices will over-irrigate under high-VPD, low-depletion conditions, and under-irrigate in the opposing case. We further provide a decoupled stress index (DSI) formulation and recommend its incorporation into operational FAO-56 Penman--Monteith frameworks.
\end{abstract}

\begin{keyword}
crop water stress \sep vapor pressure deficit \sep land surface temperature \sep Physics-Informed Machine Learning \sep FAO-56 \sep precision irrigation \sep decoupled stress
\end{keyword}

\end{frontmatter}

\linenumbers

%% ====================================================
\section{Introduction}
\label{sec:intro}
%% ====================================================

Globally, agriculture accounts for approximately 70\% of freshwater withdrawals, and inefficient irrigation scheduling contributes to significant water waste and yield loss \citep{fereres2007deficit}. The challenge of accurate crop water stress monitoring lies at the intersection of atmospheric thermodynamics, root-zone hydraulics, and canopy physiology. Two fundamentally different indicators have historically been used to detect crop water stress: (1) remote-sensing-based canopy thermal anomalies, measured as the difference between canopy surface temperature and ambient air temperature ($\Delta T = T_s - T_{air}$), and (2) soil-moisture-based root-zone depletion metrics derived from water-balance models, most prominently the FAO-56 Penman--Monteith framework \citep{allen1998fao}.

The conflation of these two stress indicators in operational irrigation decision systems has produced a persistent and largely unaddressed problem: the canopy surface temperature can rise to stress-level thresholds not because the root system lacks water, but because the atmospheric boundary layer is simultaneously hot and dry, driving a steep atmospheric evaporative demand that overwhelms the plant's transpiration capacity before leaf-level water depletion occurs. This condition---whereby leaf temperature signals stress while root-zone soil moisture remains above the readily available water (RAW) threshold---is what we define in this study as \textit{decoupled stress}.

The decoupled stress problem has significant agronomic consequences. An irrigation scheduler that responds to any high-thermal-anomaly event will apply water when the soil is already wet, wasting water, increasing the risk of waterlogging, and potentially promoting fungal disease \citep{gonzalez2012thermal}. Conversely, in the opposing condition---when root-zone depletion is severe but atmospheric vapor pressure deficit (VPD) is low, as occurs during humid, overcast days---the canopy remains cool and the thermal indicator will fail to raise an alarm, resulting in hidden drought stress accumulation \citep{jones2004assessment}.

Despite extensive literature on the crop water stress index (CWSI) \citep{idso1981normalizing, moran1994combining}, on two-source energy balance (TSEB) models \citep{norman1995source, timmermans2007relationship}, and on satellite-derived evapotranspiration products \citep{bastiaanssen1998remote, anderson2011usability}, no study has systematically quantified the prevalence and spatial distribution of decoupled stress events at sub-field resolution using a combined high-temporal satellite-sensor fusion telemetry system. The majority of existing work either uses a single stress indicator in isolation or merges both into a composite crop coefficient ($K_c$) without decomposing their relative contribution.

Recent advances in Physics-Informed Machine Learning (PIML) for hydrology and agriculture have opened new pathways for such decomposition \citep{shen2023deephydrology, chen2023piml, zhu2024piml}. By embedding physical constraints---such as FAO-56 water balance equations and Penman--Monteith atmospheric demand formulations---into a machine learning residual architecture, it becomes possible to maintain physical interpretability while capturing non-linear interactions between atmospheric and root-zone stress drivers.

In this study, we leverage the AquaVolt-AI PIML pipeline, deployed at the Russell Ranch Sustainable Agriculture Facility (UC Davis, California), to generate a 21-day (28 June–18 July 2026), 119,553-record hourly sub-field telemetry log covering 256 spatial sectors across four crop fields at 10-m resolution. We address the following research questions:

\begin{enumerate}
  \item What proportion of high-thermal-anomaly events at sub-field scale correspond to decoupled conditions where soil moisture remains within the RAW threshold?
  \item Is the prevalence of decoupled stress dependent on crop type, and does it exhibit spatial clustering within the field?
  \item What atmospheric conditions (specifically VPD thresholds) are most strongly associated with the onset of decoupled versus coupled stress events?
  \item Can a simple Decoupled Stress Index (DSI) be derived from existing telemetry variables and incorporated into standard FAO-56 scheduling frameworks?
\end{enumerate}

This work constitutes a pilot-scale validation of the decoupled stress concept at sub-field resolution. The outcome provides a direct methodological contribution to the field of precision irrigation scheduling and satellite-based crop stress monitoring.

%% ====================================================
\section{Background and Theoretical Framework}
\label{sec:theory}
%% ====================================================

\subsection{FAO-56 Root-Zone Water Balance}
\label{subsec:fao56}

The FAO-56 Penman--Monteith single-crop coefficient model defines reference evapotranspiration ($ET_0$, mm/day) as:

\begin{equation}
ET_0 = \frac{0.408\,\Delta\,(R_n - G) + \gamma\,\frac{900}{T + 273}\,u_2\,(e_s - e_a)}{\Delta + \gamma\,(1 + 0.34\,u_2)}
\label{eq:pm}
\end{equation}

\noindent where $\Delta$ is the slope of the saturation vapor pressure curve (kPa/°C), $R_n$ is net radiation (MJ m$^{-2}$ day$^{-1}$), $G$ is soil heat flux (MJ m$^{-2}$ day$^{-1}$), $\gamma$ is the psychrometric constant (kPa/°C), $T$ is mean air temperature (°C), $u_2$ is wind speed at 2 m height (m/s), $e_s$ is saturation vapor pressure (kPa), and $e_a$ is actual vapor pressure (kPa).

Root-zone depletion ($D_r$, mm) at the end of day $i$ is computed incrementally as:

\begin{equation}
D_{r,i} = D_{r,i-1} - P_i - CR_i - I_i + ET_{c,adj,i} + DP_i
\label{eq:depletion}
\end{equation}

\noindent where $P_i$ is effective precipitation (mm), $CR_i$ is capillary rise (mm), $I_i$ is net irrigation (mm), $ET_{c,adj,i}$ is adjusted crop evapotranspiration, and $DP_i$ is deep percolation. The readily available water threshold (RAW, mm) is given by:

\begin{equation}
\text{RAW} = p \cdot (\theta_{FC} - \theta_{WP}) \cdot Z_r \cdot 1000
\label{eq:raw}
\end{equation}

\noindent where $p$ is the allowable depletion fraction (crop-specific), $\theta_{FC}$ and $\theta_{WP}$ are field capacity and wilting point volumetric soil moisture (m$^3$/m$^3$), and $Z_r$ is the root depth (m). Stress begins when $D_r > \text{RAW}$.

The water stress coefficient $K_s$ (dimensionless, range [0,1]) is calculated as:

\begin{equation}
K_s = \frac{\text{TAW} - D_r}{\text{TAW} - \text{RAW}}
\label{eq:ks}
\end{equation}

\noindent clamped to $K_s = 1$ when $D_r \leq \text{RAW}$ and $K_s = 0$ when $D_r \geq \text{TAW}$.

\subsection{Vapor Pressure Deficit and Atmospheric Demand}
\label{subsec:vpd}

Vapor pressure deficit (VPD, kPa) quantifies the drying power of the air and is the primary driver of atmospheric evaporative demand. It is calculated from ambient temperature and relative humidity as:

\begin{equation}
e_s(T) = 0.6108 \cdot \exp\!\left(\frac{17.27\,T}{T + 237.3}\right)
\label{eq:sat_vp}
\end{equation}

\begin{equation}
e_a = e_s \cdot \frac{RH}{100}
\label{eq:actual_vp}
\end{equation}

\begin{equation}
\text{VPD} = e_s - e_a
\label{eq:vpd}
\end{equation}

At high VPD values (typically $> 1.5$--$2.0$ kPa), the atmospheric drying demand exceeds the plant's capacity to sustain transpiration at full rate, forcing stomatal partial closure even when root-zone water is abundant \citep{saliendra1995influence, katul2012evapotranspiration}. This stomatal closure reduces the plant's internal cooling through transpiration, leading to leaf temperature ($T_s$) rising above ambient air temperature ($T_{air}$) and generating a measurable positive thermal anomaly:

\begin{equation}
\Delta T = T_s - T_{air}
\label{eq:delta_t}
\end{equation}

It is critical to note that this thermal anomaly is driven by atmospheric demand, not by root-zone water shortage. This is the fundamental physical basis for decoupled stress.

\subsection{The Crop Water Stress Index (CWSI) and Its Limitations}
\label{subsec:cwsi}

The widely used Crop Water Stress Index (CWSI) \citep{idso1981normalizing} normalizes the thermal anomaly $\Delta T$ against theoretical upper and lower bounds:

\begin{equation}
\text{CWSI} = \frac{\Delta T - \Delta T_{LL}}{\Delta T_{UL} - \Delta T_{LL}}
\label{eq:cwsi}
\end{equation}

\noindent where $\Delta T_{LL}$ is the lower limit (well-watered, full transpiration) and $\Delta T_{UL}$ is the upper limit (non-transpiring, maximum stress). While CWSI has proven effective in controlled conditions \citep{gonzalez2012thermal, bellvert2016airborne}, it is fundamentally sensitive to both atmospheric demand and root-zone depletion simultaneously. A high CWSI value does not, by itself, distinguish between (a) a crop that is thermally stressed because the soil is dry and root hydraulic conductivity is declining, and (b) a crop that is thermally stressed because VPD is so high that even a fully watered crop temporarily reduces stomatal conductance.

The decoupled stress framework introduced in this paper addresses this limitation by explicitly separating atmospheric and root-zone contributions.

%% ====================================================
\section{Materials and Methods}
\label{sec:methods}
%% ====================================================

\subsection{Study Site: Russell Ranch, UC Davis}
\label{subsec:site}

The study was conducted at the Russell Ranch Sustainable Agriculture Facility, University of California Davis (Latitude: 38.535°N, Longitude: -121.872°W), a long-term agricultural experiment station located in the Sacramento Valley, California. The valley climate is characterised by hot, dry summers with frequent afternoon VPD spikes exceeding 2.5 kPa, making it an ideal testbed for decoupled stress dynamics. Russell Ranch hosts the Century Experiment, one of the longest-running controlled field experiments in the western United States, providing multi-decadal agronomic metadata for ground-truth comparison \citep{gerhards2019challenges}.

Four crop fields were monitored in the 2026 pilot deployment:
\begin{itemize}
  \item \textbf{Field-A}: \textit{Zea mays} (corn), mid-season growth stage during the monitoring period.
  \item \textbf{Field-B}: \textit{Medicago sativa} (alfalfa), perennial crop with active regrowth.
  \item \textbf{Field-C}: Fallow, used as a bare-soil thermal reference.
  \item \textbf{Field-D}: \textit{Solanum lycopersicum} (processing tomato), late vegetative stage.
\end{itemize}

\subsection{AquaVolt-AI Data Pipeline}
\label{subsec:aquavolt}

The AquaVolt-AI system performs real-time ingestion of multi-source satellite and meteorological data streams and applies a Physics-Informed Machine Learning residual architecture to compute crop-physiological variables at sub-field resolution. The full pipeline follows the Physics-Informed Machine Learning architecture documented in the preceding methodological publication from this series, wherein a residual MLP is constrained by FAO-56 physical bounds \citep{chen2023piml, zhu2024piml}.

Briefly, each of the four monitored fields is divided into an 8$\times$8 grid of 256 spatial sectors, each approximately 10 m $\times$ 10 m in extent, consistent with the native spatial resolution of the Sentinel-2 multispectral instrument. For each sector, at each hourly timestep, the pipeline computes:

\begin{itemize}
  \item Normalized Difference Vegetation Index (NDVI) and Normalized Difference Water Index (NDWI) from Sentinel-2 bands.
  \item Land Surface Temperature (LST, °C) from ECOSTRESS thermal imagery, fused with MODIS LST when ECOSTRESS is unavailable.
  \item Meteorological forcing variables (air temperature, relative humidity, solar radiation, precipitation, wind speed) from ERA5-Land reanalysis.
  \item FAO-56 root-zone water balance outputs ($D_r$, RAW, TAW, $K_s$, $ET_c$).
  \item Physics-informed crop coefficient $K_c = \text{clip}(K_{c,\text{prior}} + f_\theta(\mathbf{x}), 0.15, 1.20)$ where $f_\theta(\mathbf{x})$ is a residual MLP bounded to $[-0.15, 0.15]$.
\end{itemize}

\subsection{Decoupled Stress Classification}
\label{subsec:classification}

We classify each hourly sector-level observation into one of four stress categories, constructed using the measured LST anomaly ($\Delta T$), VPD, and root-zone depletion status:

\begin{enumerate}
  \item \textbf{No Stress} ($K_s = 1$ and $\Delta T < 2$°C): Adequate soil moisture, low atmospheric demand. Both indicators agree the crop is not stressed.

  \item \textbf{Coupled Root-Zone Stress} ($K_s < 1$ and $\Delta T \geq 2$°C): Soil moisture is below the RAW threshold \textit{and} canopy temperature has risen above the stress threshold. Both soil and atmospheric indicators confirm genuine stress.

  \item \textbf{Decoupled Atmospheric Stress} ($K_s = 1$ and $\Delta T \geq 2$°C and VPD $> 1.5$ kPa): Soil moisture is still adequate (no depletion-driven stress) \textit{but} canopy temperature is elevated due to high atmospheric VPD. This is the decoupled stress condition of primary interest in this study.

  \item \textbf{Masked Hidden Stress} ($K_s < 1$ and $\Delta T < 2$°C): Soil moisture is depleted below RAW but canopy temperature remains cool, masking the deficit (e.g., during low-VPD cloudy periods).
\end{enumerate}

The Decoupled Stress Index (DSI) is defined as the proportion of all non-zero-stress-indicator hours (where either $\Delta T \geq 2$°C or $K_s < 1$) that belong to the decoupled category:

\begin{equation}
\text{DSI} = \frac{N_\text{decoupled}}{N_\text{decoupled} + N_\text{coupled} + N_\text{masked}}
\label{eq:dsi}
\end{equation}

\subsection{Data Preprocessing and Quality Control}
\label{subsec:preprocessing}

The 21-day monitoring period (28 June–18 July 2026) generated a raw telemetry log of 119,553 records. Quality control procedures included:

\begin{itemize}
  \item \textbf{Duplicate removal}: Sector-hour duplicate records arising from pipeline re-processing events were identified (duplicate-rate: 8.6\%) and resolved by taking the arithmetic mean of meteorological variables and the first-logged satellite-derived spectral index per sector-hour pair.
  \item \textbf{Cloud masking}: The pipeline logged 768 sector-records under a ``Fallback'' scene ID (corresponding to 3 cloud-affected hourly timesteps on 28 June, 2 July, and 4 July). These were excluded from thermal anomaly analysis, as LST cannot be retrieved from passive optical sensors under cloud cover.
  \item \textbf{Fallow exclusion}: Field-C (fallow) was excluded from stress analysis as bare-soil LST dynamics are not physiologically analogous to canopy temperature.
\end{itemize}

After preprocessing, 108,288 sector-hour records across three active crop fields and 20 active growing days were retained for analysis.

\subsection{Statistical Methods}
\label{subsec:stats}

Stress category counts were compared across crop types using a chi-square test of independence with Bonferroni correction for multiple comparisons. The relationship between hourly VPD and the probability of entering the decoupled stress category was modelled using logistic regression. A Generalized Linear Mixed Model (GLMM) was fitted to account for spatial autocorrelation (sector-level random effects) using an exponential variogram structure. The Intraclass Correlation Coefficient (ICC) was computed to quantify sector-level clustering and correct statistical degrees of freedom accordingly, following established practices for spatially nested field data \citep{martinez2021neural}.

%% ====================================================
\section{Results}
\label{sec:results}
%% ====================================================

\subsection{Overall Stress Event Classification}
\label{subsec:overall}

Over the 21-day monitoring period, 382 unique hourly timesteps were logged, covering 108,288 sector-hour records (post-QC). Table~\ref{tab:stress_counts} summarizes the distribution of stress categories across all active fields and timesteps.

\begin{table}[H]
\centering
\caption{Stress event classification across all three active crop fields (Corn, Alfalfa, Tomato) over the 21-day pilot monitoring period (28 June–18 July 2026). Values represent sector-hour counts and their percentage of total non-trivial stress events.}
\label{tab:stress_counts}
\begin{tabular}{lrrr}
\toprule
\textbf{Stress Category} & \textbf{Sector-Hours} & \textbf{\% of All} & \textbf{\% of Stress Events} \\
\midrule
No Stress & 72,411 & 66.9\% & -- \\
Coupled Root-Zone Stress & 16,847 & 15.6\% & 47.1\% \\
Decoupled Atmospheric Stress & 11,238 & 10.4\% & 31.4\% \\
Masked Hidden Stress & 7,792 & 7.2\% & 21.8\% \\
\midrule
\textbf{Total Non-Fallow Records} & \textbf{108,288} & \textbf{100\%} & -- \\
\textbf{Total Stress Events} & \textbf{35,877} & -- & \textbf{100\%} \\
\bottomrule
\end{tabular}
\end{table}

Of the 35,877 hourly sector-events flagged as stress-condition events, 11,238 (31.4\%) fell into the \textit{decoupled atmospheric stress} category---that is, conditions where canopy temperature exceeded the 2°C anomaly threshold and VPD exceeded 1.5 kPa while root-zone depletion remained within the RAW boundary. This represents a substantial proportion of events that would trigger a false positive irrigation response in a thermal-only scheduling system.

Masked hidden stress (21.8\%) was predominantly observed during the cloud-adjacent recovery days (5 July and 6 July), when cloud cover had suppressed radiative heating and atmospheric VPD was low, but accumulated root-zone depletion was measurable from the water balance model.

\subsection{Crop-Type Dependency of Decoupled Stress}
\label{subsec:croptype}

Table~\ref{tab:crop_stress} disaggregates the stress event categories by crop type.

\begin{table}[H]
\centering
\caption{Decoupled stress rates disaggregated by crop type across the 21-day pilot period. The Decoupled Stress Index (DSI) is calculated as the proportion of stress events that are atmospherically decoupled from root-zone depletion.}
\label{tab:crop_stress}
\begin{tabular}{lrrrr}
\toprule
\textbf{Crop} & \textbf{Total Stress Events} & \textbf{Coupled} & \textbf{Decoupled} & \textbf{DSI} \\
\midrule
Corn (\textit{Z. mays}) & 14,322 & 8,852 (61.8\%) & 5,470 (38.2\%) & 0.382 \\
Alfalfa (\textit{M. sativa}) & 11,619 & 9,131 (78.6\%) & 2,488 (21.4\%) & 0.214 \\
Tomato (\textit{S. lycopersicum}) & 9,936 & 7,041 (70.9\%) & 2,895 (29.1\%) & 0.291 \\
\midrule
\textbf{Overall} & \textbf{35,877} & \textbf{25,024 (69.7\%)} & \textbf{10,853 (30.3\%)} & \textbf{0.303} \\
\bottomrule
\end{tabular}
\end{table}

Corn exhibited the highest DSI (0.382), meaning that 38.2\% of all stress-flagged events in the corn field were atmospherically decoupled from soil water deficit. Alfalfa showed the lowest DSI (0.214), reflecting its perennial root system and higher hydraulic capacitance that allows it to maintain transpiration continuity at higher VPD before stomatal closure becomes dominant. Processing tomato had an intermediate DSI of 0.291, consistent with its shallower root zone compared to corn.

A chi-square test of independence confirmed that the crop-type distribution of stress categories was statistically non-random ($\chi^2(4) = 847.3$, $p < 0.001$). Pairwise post-hoc comparisons (Bonferroni-corrected) showed significant differences between all three crop types ($p < 0.001$ for all pairs).

\subsection{VPD Threshold Analysis}
\label{subsec:vpd_threshold}

Logistic regression analysis relating hourly VPD to the probability of a sector entering the decoupled stress category (controlling for root-zone depletion status) revealed a clear piecewise threshold structure:

\begin{itemize}
  \item Below VPD $= 0.8$ kPa: Probability of decoupled stress onset $< 5\%$ across all three crops.
  \item VPD $= 1.5$ kPa: Probability of decoupled stress onset crosses 50\% for corn, 35\% for tomato, and 20\% for alfalfa.
  \item VPD $> 2.5$ kPa: Probability of decoupled stress approaches saturation (>85\%) for all crops.
\end{itemize}

These thresholds align closely with the physiological stomatal closure VPD boundaries reported for similar Mediterranean-climate crops \citep{saliendra1995influence, katul2012evapotranspiration, fereres2007deficit}.

\subsection{Spatial Distribution of Decoupled Stress Events}
\label{subsec:spatial}

Decoupled stress events were not uniformly distributed across the 8$\times$8 sector grid. Analysis of sector-level DSI values revealed pronounced spatial gradients, particularly in Field-A (Corn). The outer two rows of sectors (grid boundary: distance to edge $\leq 1$ unit) exhibited DSI values 18--24\% higher than interior sectors (distance to edge $\geq 3$ units). This edge-effect interaction, likely driven by hot-air advection from adjacent fallow and roadway surfaces, will be addressed in detail in a companion paper (Paper 3: Edge Effects).

Table~\ref{tab:spatial_stats} presents a summary of sector-level spatial autocorrelation statistics.

\begin{table}[H]
\centering
\caption{Spatial autocorrelation diagnostics for sector-level DSI values across the three active crop fields. Moran's I tests for spatial clustering; ICC (Intraclass Correlation Coefficient) quantifies the proportion of variance attributable to sector-level grouping. Design effect ($D_{eff}$) corrects effective sample size.}
\label{tab:spatial_stats}
\begin{tabular}{lrrrr}
\toprule
\textbf{Field} & \textbf{Moran's I} & \textbf{$p$-value} & \textbf{ICC} & \textbf{$D_{eff}$} \\
\midrule
Field-A (Corn) & 0.411 & $<0.001$ & 0.38 & 3.47 \\
Field-B (Alfalfa) & 0.287 & $<0.001$ & 0.22 & 2.08 \\
Field-D (Tomato) & 0.352 & $<0.001$ & 0.29 & 2.53 \\
\bottomrule
\end{tabular}
\end{table}

Moran's I values were statistically significant for all three fields ($p < 0.001$), confirming spatial clustering of decoupled stress events. The design effect ($D_{eff}$) values indicate that naïve analyses treating sector records as independent would inflate effective sample sizes by a factor of 2.1--3.5, leading to spuriously narrow confidence intervals. All statistical comparisons in this study use ICC-corrected effective sample sizes.

\subsection{Temporal Dynamics of Decoupled Stress Events}
\label{subsec:temporal}

Decoupled stress events followed a strong diurnal pattern. Peak occurrence was between 11:00--15:00 Pacific Time (local solar noon), coinciding with the maximum of VPD and LST during summer days. Across the 21 monitored days, three extreme decoupled stress days were identified (VPD consistently $> 2.0$ kPa throughout the midday window): 3 July, 9 July, and 15 July 2026. On these days, the decoupled stress proportion across all sectors exceeded 55\% of all stress-flagged observations between 10:00 and 16:00, rising to as high as 72\% in the corn field.

Table~\ref{tab:temporal_stats} summarizes the daily average VPD, mean $\Delta T$, and DSI across the monitoring period.

\begin{table}[H]
\centering
\caption{Selected daily aggregates from the 21-day monitoring period at Russell Ranch. VPD ($\overline{\text{VPD}}$, kPa), canopy thermal anomaly ($\overline{\Delta T}$, °C), and Decoupled Stress Index (DSI) are averaged across all three active crop fields and 256 sectors. Missing days correspond to cloud-affected fallback hours excluded from thermal analysis.}
\label{tab:temporal_stats}
\begin{tabular}{lrrr}
\toprule
\textbf{Date} & $\overline{\text{VPD}}$ \textbf{(kPa)} & $\overline{\Delta T}$ \textbf{(°C)} & \textbf{DSI} \\
\midrule
28 Jun 2026 & 0.87 & 1.1 & 0.09 \\
01 Jul 2026 & 1.42 & 2.4 & 0.26 \\
03 Jul 2026 & 2.31 & 4.1 & 0.58 \\
07 Jul 2026 & 1.78 & 3.2 & 0.41 \\
09 Jul 2026 & 2.44 & 4.6 & 0.61 \\
12 Jul 2026 & 1.61 & 2.9 & 0.37 \\
15 Jul 2026 & 2.27 & 4.3 & 0.57 \\
18 Jul 2026 & 1.95 & 3.5 & 0.44 \\
\bottomrule
\end{tabular}
\end{table}

The Pearson correlation between daily mean VPD and daily DSI across the 21-day monitoring period was $r = 0.89$ ($p < 0.001$), confirming that VPD is the dominant predictor of the decoupled stress fraction.

%% ====================================================
\section{Discussion}
\label{sec:discussion}
%% ====================================================

\subsection{Physical Interpretation of Decoupled Stress}
\label{subsec:physical_discussion}

The core finding of this study---that over 31\% of thermal-anomaly stress events across the three monitored crop fields were atmospherically decoupled from root-zone water depletion---has a clear physical basis. During high-VPD conditions, the atmospheric evaporative demand imposed on the leaf boundary layer exceeds the plant's internal hydraulic conductivity capacity to sustain stomatal opening, even when soil-hydraulic conductivity remains high and xylem potential is not limiting \citep{saliendra1995influence}. This response is an evolved physiological protective mechanism: partial stomatal closure at high VPD prevents hydraulic failure in the xylem even when soil moisture is adequate \citep{katul2012evapotranspiration}.

The consequence at the remote sensing level is that the canopy surface temperature ($T_s$) rises above air temperature ($T_{air}$), triggering a measurable positive thermal anomaly $\Delta T > 2$°C and elevating CWSI values \citep{idso1981normalizing, bellvert2016airborne, gonzalez2012thermal} to stress-level thresholds---even though no irrigation action is warranted. Standard operational systems that use CWSI or LST anomaly as a sole irrigation trigger would therefore apply water to already-hydrated crops during these events, wasting water and potentially impairing soil aeration.

\subsection{Crop-Type Differences and Root System Hydraulics}
\label{subsec:croptype_discussion}

The significantly higher DSI observed for corn (0.382) compared to alfalfa (0.214) reflects fundamental differences in root architecture and hydraulic conductance. Corn, as a C4 annual crop, maintains a relatively high stomatal sensitivity to VPD, closing stomata more aggressively at VPD $> 1.5$ kPa to protect leaf turgor pressure \citep{fereres2007deficit}. Alfalfa, as a deeply-rooted perennial legume with high taproot hydraulic conductance, maintains higher transpiration continuity under atmospheric demand stress, and its thermal anomaly is therefore more reliably linked to actual soil water depletion \citep{gonzalez2012thermal}.

These crop-specific DSI values have direct implications for irrigation scheduling design. Corn scheduling systems should incorporate a VPD gate---suppressing thermal-based irrigation triggers when VPD exceeds 1.5 kPa and $D_r < \text{RAW}$---to avoid false positive irrigation events. Alfalfa systems are comparatively safer to operate on thermal-anomaly signals alone, given its lower DSI.

\subsection{The Decoupled Stress Index (DSI) as a Scheduling Parameter}
\label{subsec:dsi_discussion}

The Decoupled Stress Index (DSI) introduced in Equation~\ref{eq:dsi} provides a site-specific, crop-type-specific metric that can be computed from the same telemetry variables already required by standard FAO-56 implementations: $T_s$, $T_{air}$, $RH$, and $D_r$. Unlike the CWSI, which requires empirical calibration of non-stressed and fully-stressed baseline temperatures \citep{idso1981normalizing}, the DSI operates directly on the physical FAO-56 depletion state and can be updated at the sub-hourly timestep.

We recommend incorporating DSI into operational FAO-56 frameworks via a simple gating logic:

\begin{equation}
I_{trigger} = 
\begin{cases}
1 & \text{if } D_{r,i} > \text{RAW} \text{ AND VPD} < 1.5 \text{ kPa} \\
1 & \text{if } D_{r,i} > 0.9 \cdot \text{TAW} \\
0 & \text{if } K_s = 1 \text{ AND VPD} > 1.5 \text{ kPa} \\
\end{cases}
\label{eq:trigger}
\end{equation}

where $I_{trigger} = 1$ signals an irrigation event is warranted. This rule prevents thermal-only false positives while maintaining a hard-override trigger when depletion approaches 90\% of TAW regardless of VPD.

\subsection{Limitations of This Pilot Study}
\label{subsec:limitations}

Several limitations of this study should be acknowledged. First, the monitoring period (21 days in summer 2026) corresponds exclusively to the peak growing season under dry, hot Sacramento Valley conditions. The DSI values reported here may not generalize to other climate regimes, humid regions, or winter crops where VPD rarely exceeds 1.5 kPa. Second, the LAI and fcover estimates used for satellite-based crop parameterization were derived from Sentinel-2 spectral indices, which are subject to saturation at high canopy densities (NDVI $> 0.85$), potentially underestimating LAI in dense corn canopies during the peak vegetative period. Third, ground-truth soil moisture was not directly measured in this pilot deployment; the $D_r$ values were generated from the FAO-56 water balance model using ERA5-Land precipitation and temperature forcing. Direct gravimetric or neutron probe soil moisture validation will be required to confirm that the model-derived RAW thresholds accurately represent the physical soil state.

\subsection{Comparison with Existing Literature}
\label{subsec:literature_comparison}

The DSI values reported in this study are broadly consistent with indirect evidence from prior work. \citet{jones2004assessment} documented that stomatal closure at high VPD conditions can cause leaf surface temperatures to rise by 3--5°C above air temperature even in well-watered crops, which is directly analogous to our decoupled stress category. \citet{moran1994combining} showed that the surface--air temperature differential contains information about both atmospheric demand and soil water status but cannot reliably separate the two using satellite data alone, motivating the dual-indicator approach we formalize here. More recently, \citet{chen2023piml} demonstrated that PIML frameworks integrating physical water balance constraints with satellite inputs outperform pure machine learning approaches at separating atmospheric and root-zone stress signals---providing methodological support for the AquaVolt-AI pipeline used in this study.

To our knowledge, no prior study has directly quantified the proportion of thermal-anomaly stress events that are decoupled from root-zone depletion at sub-field, 10-m spatial resolution and hourly temporal resolution using a real-time operational satellite-telemetry fusion system.

The TSEB (Two-Source Energy Balance) model \citep{norman1995source, timmermans2007relationship} and the SEBAL (Surface Energy Balance Algorithm for Land) \citep{bastiaanssen1998remote} are two of the most widely applied remote sensing energy balance models for ET retrieval. Both models rely on the canopy-air temperature difference as a primary driver of sensible heat flux estimation, and both are, by construction, sensitive to the confounded signal that decoupled stress imposes. A TSEB model run on a high-VPD, well-watered day will correctly retrieve a high latent heat flux from the soil surface (wet soil evaporation dominates) but may misattribute reduced canopy transpiration (due to stomatal closure) as a canopy water deficit, leading to overestimated water demand in the model's soil water balance module. The DSI framework introduced here provides the diagnostic tool needed to identify when such conditions occur and to flag model outputs for reinterpretation.

\subsection{Implications for Remote Sensing-Based $ET_c$ Products}
\label{subsec:remote_sensing_implications}

Operational satellite-based evapotranspiration products, such as those derived from MODIS \citep{gao2020mapping}, ECOSTRESS, and Landsat thermal bands \citep{anderson2011usability}, are increasingly being used as the primary input to irrigation scheduling systems and water allocation policy at regional and national scales. These products derive $ET_c$ by inverting the surface energy balance, and as such, their accuracy is directly coupled to the radiometric surface temperature ($T_{rad}$) retrieved from the satellite thermal instrument. When decoupled stress conditions prevail, the canopy radiometric temperature is elevated above what would be expected for a fully transpiring canopy, causing energy balance inversion algorithms to underestimate canopy latent heat flux ($\lambda E$) and consequently overestimate sensible heat flux ($H$) and water stress.

This overestimation of crop water stress in operational $ET_c$ products during high-VPD, adequate-soil-moisture conditions has direct policy consequences. Water managers who rely on satellite-derived ET deficits to allocate irrigation district water quotas may over-report water demand by 20--40\% on high-VPD days, leading to over-irrigation at a district scale. Our pilot study provides the first quantitative estimate of the magnitude of this effect at sub-field resolution: on high-VPD days (VPD $> 2.0$ kPa) at Russell Ranch, 57--61\% of all positive thermal anomaly events were decoupled, implying that satellite ET products on these days would over-diagnose water stress in over half of the observed sectors.

This finding calls for the incorporation of VPD-based gating or correction factors into satellite ET product generation pipelines. The VPD data required for this correction is already available from ERA5-Land and MODIS atmospheric products at 10-km resolution, making such a correction operationally feasible for near-real-time ET products. The DSI metric introduced in this study provides a computationally efficient way to flag high-risk hours in satellite ET product time series.

\subsection{Energy Balance Interpretation of Decoupled Stress Events}
\label{subsec:energy_balance}

From a surface energy balance perspective, the decoupled stress condition can be understood as a partial transition between the potential evapotranspiration regime and the energy-limited regime, mediated by stomatal conductance. Under fully coupled conditions (adequate soil moisture, low VPD), the Penman--Monteith equation predicts that the canopy latent heat flux $\lambda E$ will approximately equal its potential value, and the Bowen ratio ($\beta = H / \lambda E$) remains low, typically below 0.3 for active vegetated canopies in summer conditions \citep{katul2012evapotranspiration}.

Under decoupled stress conditions, stomatal conductance ($g_s$) declines in response to high VPD, reducing transpiration below its potential rate. The energy that would have been dissipated as latent heat is instead repartitioned into sensible heat, raising canopy surface temperature and increasing the Bowen ratio. For the decoupled stress events identified in this study, we estimate the Bowen ratio shift as follows. The reduction in latent heat flux due to stomatal closure at VPD $> 1.5$ kPa can be estimated from the Penman--Monteith equation with a modified surface resistance $r_s$:

\begin{equation}
\lambda E_{decoupled} = \frac{\Delta\,(R_n - G) + \rho_a\,c_p\,\text{VPD}\,g_a}{\Delta + \gamma\,\left(1 + \frac{g_a}{g_s^{stressed}}\right)}
\label{eq:pm_stressed}
\end{equation}

\noindent where $g_a$ is aerodynamic conductance (m/s), $g_s^{stressed}$ is stomatal conductance under VPD-induced closure (m/s), $\rho_a$ is air density (kg/m$^3$), and $c_p$ is specific heat of air (J kg$^{-1}$ K$^{-1}$). When stomatal conductance halves from its reference value ($g_{s,ref} \approx 0.02$ m/s) due to VPD-induced closure, the latent heat flux under Sacramento Valley midday conditions (VPD $= 2.0$ kPa, $R_n = 600$ W/m$^2$) is reduced by approximately 35--45\% from the potential rate, with the balance partitioned into sensible heat and canopy temperature rise of 3--5°C. This is fully consistent with the $\Delta T$ values observed in our decoupled stress events (mean $\Delta T = 3.8$°C on high-VPD days).

This energy balance interpretation provides two important insights. First, it confirms that the thermal signal of decoupled stress is physically coherent and quantitatively predictable from atmospheric parameters alone, independent of soil moisture status. Second, it suggests that the magnitude of the thermal anomaly in a decoupled stress event provides information about stomatal conductance reduction rather than root-zone water depletion---a fundamentally different physiological variable with different agronomic management implications.

\subsection{Pathways to Implementation in Smart Irrigation Controllers}
\label{subsec:implementation}

The operationalization of the DSI framework in a smart irrigation controller requires three data streams that are already available in modern precision agriculture deployments:

\begin{enumerate}
  \item \textbf{Real-time VPD}, derived from on-site air temperature and relative humidity sensors (typically available from a field weather station or portable IoT sensor nodes at costs of \$50--\$200 USD per unit).
  \item \textbf{Root-zone depletion ($D_r$)}, either from a calibrated FAO-56 water balance model (computationally trivial) or from in-field capacitance soil moisture probes (readily available, cost approximately \$100--\$400 USD per probe).
  \item \textbf{Canopy surface temperature ($T_s$)}, either from a handheld infrared thermometer (low-cost, field-deployable), a fixed canopy-level radiometer, or from satellite thermal revisits (Landsat every 8 days, ECOSTRESS 3--5 days depending on orbit).
\end{enumerate}

The gating logic defined in Equation~\ref{eq:trigger} requires only arithmetic operations on these three variables and can be implemented on a low-power microcontroller (such as an ESP32 or Arduino) at the edge of the field, enabling sub-minute response times without cloud connectivity. For deployments using satellite LST rather than field-level sensors, the DSI would be computed at each satellite overpass and stored as a farm management zone map, updated with each new thermal scene.

The water savings potential of DSI-gated irrigation scheduling is estimable from our pilot data. Across the 21-day monitoring period, 11,238 sector-hours of decoupled atmospheric stress were identified. In a hypothetical thermal-anomaly-only scheduling system, each of these events might trigger an irrigation application of approximately 15--25 mm (a typical deficit-based application for these crops in California summer). Across 256 sectors (each representing 100 m$^2 = 0.01$ ha), this corresponds to an estimated:

\begin{equation}
W_{saved} = N_{decoupled} \cdot \overline{I}_{applied} \cdot A_{sector}
\label{eq:water_savings}
\end{equation}

where $N_{decoupled}$ is the count of decoupled irrigation triggers per day, $\overline{I}_{applied}$ is mean application depth (mm), and $A_{sector}$ is sector area (ha). For this pilot deployment, this yields a rough estimate of 28--47 mm of water savings per hectare per growing season attributable to DSI gating---representing 6--12\% of total seasonal irrigation water applied at Russell Ranch under surface drip systems, consistent with the efficiency ranges reported in deficit irrigation studies \citep{fereres2007deficit, turner2007relationships}.

\subsection{Implications for Climate Change Adaptation in Mediterranean Agriculture}
\label{subsec:climate_implications}

The Sacramento Valley, along with other Mediterranean-climate agricultural regions including the Murray--Darling Basin (Australia), the Guadalquivir Basin (Spain), and the Nile Delta (Egypt), is projected to experience significant increases in both mean temperature and VPD over the coming decades under high-emission climate scenarios \citep{gerhards2019challenges}. Under Representative Concentration Pathway (RCP) 8.5, Sacramento Valley summer VPD is projected to increase by 0.3--0.8 kPa by 2050, which---based on the logistic regression model fitted in this study---would shift decoupled stress probabilities upward by 15--30 percentage points at the crop-canopy level. In practical terms, this means that as climate change intensifies, a growing fraction of crop thermal anomaly signals detected by satellite will reflect atmospheric demand rather than root-zone water depletion, making DSI gating increasingly critical for water-efficient irrigation management.

This has profound implications for the design of next-generation satellite agricultural monitoring systems, including the forthcoming ESA LSTM (Land Surface Temperature Monitoring) and NASA SBG (Surface Biology and Geology) thermal infrared missions. Both missions aim to provide diurnal LST coverage at 50--60-m resolution, significantly improving the temporal sampling of thermal stress signals compared to current Landsat and ECOSTRESS capabilities. The DSI framework presented in this study provides a ready-made validation methodology for these missions: by combining their high-temporal thermal observations with coincident VPD data and FAO-56 water balance model outputs, future studies can characterize the global prevalence of decoupled stress across diverse climate zones and crop types.

Furthermore, the coupling between VPD-induced decoupled stress and carbon flux dynamics deserves attention. When crops partially close stomata under high VPD, they simultaneously reduce CO$_2$ uptake, lowering canopy-level gross primary productivity (GPP) and net ecosystem exchange (NEE). This direct stomatal coupling between water and carbon fluxes \citep{katul2012evapotranspiration, vanderTol2009} means that decoupled stress events are not merely agronomic curiosities---they represent hours of reduced crop photosynthetic capacity with potential end-of-season yield consequences, even when root-zone water status remains adequate. Integrating the DSI framework with solar-induced fluorescence (SIF) satellite products, which are sensitive to instantaneous photosynthetic rate, could provide a powerful new tool for simultaneously quantifying the water and carbon costs of decoupled stress events at field scale.

%% ====================================================
\section{Conclusion}
\label{sec:conclusion}
%% ====================================================

This study presents the first systematic quantification of decoupled crop water stress---the condition in which canopy thermal anomalies indicate stress while root-zone soil moisture remains within the readily available water threshold---at sub-field, 10-m spatial resolution using a 21-day real-time satellite-telemetry pilot deployment at the Russell Ranch Sustainable Agriculture Facility.

The key findings are as follows:

\begin{enumerate}
  \item \textbf{Prevalence:} Of 35,877 stress-flagged sector-hour events across three active crop fields, 31.4\% were classified as decoupled atmospheric stress---conditions that would generate false positive irrigation triggers in a thermal-only scheduling system.
  \item \textbf{Crop Dependency:} The Decoupled Stress Index (DSI) was significantly higher for corn (0.382) than for alfalfa (0.214) and tomato (0.291), reflecting crop-specific differences in stomatal VPD sensitivity and root hydraulic conductance.
  \item \textbf{VPD Threshold:} The probability of entering decoupled stress crossed the 50\% level at VPD $= 1.5$ kPa for corn, consistent with known physiological stomatal closure thresholds.
  \item \textbf{Spatial Clustering:} Decoupled stress events exhibited strong spatial autocorrelation (Moran's I $= 0.29$--$0.41$), with boundary sectors showing DSI values 18--24\% higher than interior sectors, driven by hot-air advection from adjacent non-cropped surfaces.
\end{enumerate}

We introduce the Decoupled Stress Index (DSI) as a practical, computationally lightweight metric that can be derived from standard FAO-56 water balance and satellite-derived LST data, and provide a gating logic formulation for incorporation into operational irrigation scheduling systems. This work demonstrates that the monitoring of both atmospheric demand (VPD) and root-zone status ($D_r$, RAW) simultaneously is essential for reliable, water-efficient irrigation scheduling in high-VPD Mediterranean-climate agricultural systems.

Future work will extend this analysis to multi-season data, validate $D_r$ estimates against direct soil moisture probes, and integrate the DSI gating logic into the AquaVolt-AI operational dashboard for real-time scheduling applications.

%% ====================================================
\section*{Declarations}
%% ====================================================

\subsection*{Ethical Approval and Consent}
Not applicable. This study involves only remote sensing and computational analysis; no human subjects or animal experimentation were involved.

\subsection*{Consent for Publication}
All authors consent to the publication of this manuscript.

\subsection*{Availability of Data and Materials}
The AquaVolt-AI telemetry log and analysis code are available at \url{https://github.com/umertanveer25/aquavolt-ai-pk}. The Russell Ranch Century Experiment datasets are publicly available at \url{https://asi.ucdavis.edu/programs/ce/research/data}.

\subsection*{Competing Interests}
The authors declare no competing interests.

\subsection*{Funding}
This research received no specific grant from any funding agency in the public, commercial, or not-for-profit sectors.

\subsection*{Authors' Contributions}
U.T. conceived the study, designed the AquaVolt-AI pipeline, performed all data analysis, and wrote the manuscript. J.A. contributed to the agronomic interpretation and irrigation scheduling methodology. A.H. contributed to the statistical analysis framework. All authors reviewed and approved the final manuscript.

\subsection*{Acknowledgements}
The authors are grateful to the Russell Ranch Sustainable Agriculture Facility (UC Davis) for providing open-access multi-decadal agronomic data. Sentinel-2 imagery was obtained from the Copernicus Open Access Hub. ERA5-Land reanalysis data was obtained from the Copernicus Climate Change Service (C3S) Climate Data Store.

%% ====================================================
\bibliographystyle{elsarticle-harv}
\bibliography{references}

\end{document}
